\section{Грассманова алгебра во всём обличье}

Пусть $V$ --- модуль над кольцом~$K$ (не обязательно свободный). Рассмотрим
\[
	V^{\otimes n} :=
	\underbrace{%
		V \oco V
	}_{\text{$n$ раз}}
	.
\]
Пусть $V^{\otimes 0} := K$, $V^{\otimes 1} := V$. Тогда определим
\[
	TV :=
	\bigoplus_{n \geq 0}
		V^{\otimes n}
	.
\]
Это универсальная ассоциативная $K$-алгебра, порождённая элементами $v \in V$. Её элементы --- это мономы длины~$n$
\[
	v_1 \oco v_n \in V^{\otimes n}
\]
при $n \geq 0$ и их всевозможные $K$-линейные комбинации.

Рассмотрим идеал~$\Jskew \leq TV$, порождённый мономами вида~$v \otimes v$, где $v \in V$, то есть
\[
	\Jskew =
	\left< v \otimes v \right>_{TV} =
	\left< \cdots \otimes v \otimes v \otimes \cdots \right>,
\]
где многоточиями обозначены произвольные мономы из~$TV$. Иными словами, $\Jskew$ порождён всеми $TV$-линейными комбинациями мономов~$v \otimes v$, или, что то же самое, $\Jskew$ порождён всеми мономами, в которых есть два рядом стоящих одинаковых сомножителя.

Рассмотрим факторалгебру
\[
	\Lambda V := TV / \Jskew,
\]
в которой все элементы идеала~$\Jskew$ равны нулю. Она распадается в прямую сумму
\[
	\Lambda V =
	\bigopus_{n \geq 0}
		\Lambda^n V,
\]
где
\[
	\Lambda^n V := V^{\otimes n} / (\Jskew \cap V^{\otimes n})
	,
\]
а $\Jskew \cap V^{\otimes n}$ состоит из тензоров вида~$\cdots \otimes v \otimes v \otimes \cdots$, состоящих из~$n$ сомножителей.

Рассмотрим эпиморфизм
\begin{align*}
	\pi_{\Lambda^n} \colon
	V^{\otimes n}
	&\to
	\Lambda^n V
	\colon
	\\
	v_1 \oco v_n
	\mapsto
	[v_1 \oco v_n]_{\Lambda^n} =:
	v_1 \wcw v_n
	.
\end{align*}

Штука вида $v_1 \wcw v_n$ называется \tdef{грассмановым мономом}. 


Для любого грассманова монома $v_1 \wcw v_n \in \Lambda^n V$ выполняется следующее свойство: если где-то нашлись два рядом стоящих одинаковых сомножителя $w$, то
\[
	v_1 \wcw w \wedge w \wcw v_n = 0
	.
\]

Поглядим внимательнее на свойство $v \wedge v = 0$. Представим вектор~$v$ в виде~$v = a + b$ и подставим:
\begin{align*}
	(a + b) \wedge (a + b) &= 0,
	\\
	a \wedge a + a \wedge b + b \wedge a + b \wedge b &= 0.
\end{align*}
По определению $a \wedge a = 0$ и $b \wedge b = 0$, так что
\[
	a \wedge b + b \wedge a = 0,
\]
или, что то же самое,
\[
	a \wedge b = - b \wedge a
	.
\]
Говорят, что грассманово умножение \tdef{некоммутативно}, или \tdef{косокоммутативно}: при перемене двух соседних сомножителей местами меняется знак.

Из этого утверждения есть следствие. Пусть $\sigma \in S_n$ --- перестановка из~$n$ элементов. Тогда
\[
	v_{\sigma(1)} \wcw
	v_{\sigma(n)} =
	\sgn{\sigma}
	\cdot
	v_1 \wcw v_n
	.
\]


